% Options for packages loaded elsewhere
\PassOptionsToPackage{unicode}{hyperref}
\PassOptionsToPackage{hyphens}{url}
%
\documentclass[
  12pt,
  a4paper,
]{article}
\usepackage{amsmath,amssymb}
\usepackage{iftex}
\ifPDFTeX
  \usepackage[T1]{fontenc}
  \usepackage[utf8]{inputenc}
  \usepackage{textcomp} % provide euro and other symbols
\else % if luatex or xetex
  \usepackage{unicode-math} % this also loads fontspec
  \defaultfontfeatures{Scale=MatchLowercase}
  \defaultfontfeatures[\rmfamily]{Ligatures=TeX,Scale=1}
\fi
\usepackage{lmodern}
\ifPDFTeX\else
  % xetex/luatex font selection
  \setmainfont[]{DejaVu Serif}
  \setsansfont[]{DejaVu Sans}
\fi
% Use upquote if available, for straight quotes in verbatim environments
\IfFileExists{upquote.sty}{\usepackage{upquote}}{}
\IfFileExists{microtype.sty}{% use microtype if available
  \usepackage[]{microtype}
  \UseMicrotypeSet[protrusion]{basicmath} % disable protrusion for tt fonts
}{}
\makeatletter
\@ifundefined{KOMAClassName}{% if non-KOMA class
  \IfFileExists{parskip.sty}{%
    \usepackage{parskip}
  }{% else
    \setlength{\parindent}{0pt}
    \setlength{\parskip}{6pt plus 2pt minus 1pt}}
}{% if KOMA class
  \KOMAoptions{parskip=half}}
\makeatother
\usepackage{xcolor}
\usepackage[left=15mm,right=15mm,top=20mm,bottom=20mm]{geometry}
\usepackage{longtable,booktabs,array}
\usepackage{calc} % for calculating minipage widths
% Correct order of tables after \paragraph or \subparagraph
\usepackage{etoolbox}
\makeatletter
\patchcmd\longtable{\par}{\if@noskipsec\mbox{}\fi\par}{}{}
\makeatother
% Allow footnotes in longtable head/foot
\IfFileExists{footnotehyper.sty}{\usepackage{footnotehyper}}{\usepackage{footnote}}
\makesavenoteenv{longtable}
\usepackage{graphicx}
\makeatletter
\def\maxwidth{\ifdim\Gin@nat@width>\linewidth\linewidth\else\Gin@nat@width\fi}
\def\maxheight{\ifdim\Gin@nat@height>\textheight\textheight\else\Gin@nat@height\fi}
\makeatother
% Scale images if necessary, so that they will not overflow the page
% margins by default, and it is still possible to overwrite the defaults
% using explicit options in \includegraphics[width, height, ...]{}
\setkeys{Gin}{width=\maxwidth,height=\maxheight,keepaspectratio}
% Set default figure placement to htbp
\makeatletter
\def\fps@figure{htbp}
\makeatother
\setlength{\emergencystretch}{3em} % prevent overfull lines
\providecommand{\tightlist}{%
  \setlength{\itemsep}{0pt}\setlength{\parskip}{0pt}}
\setcounter{secnumdepth}{-\maxdimen} % remove section numbering
\usepackage{sectsty}
\allsectionsfont{\raggedright}
\usepackage{lineno}
\linenumbers
\usepackage{setspace}
\doublespacing
\usepackage{fancyhdr}
\pagestyle{fancy}\fancyhf{}
\fancyhead[L]{\small Tidal potential measured by optical clocks}
\fancyfoot[C]{\thepage}\setlength{\headheight}{14pt}
\usepackage{float}\floatplacement{figure}{H}
\setlength{\emergencystretch}{2em}
\ifLuaTeX
  \usepackage{selnolig}  % disable illegal ligatures
\fi
\IfFileExists{bookmark.sty}{\usepackage{bookmark}}{\usepackage{hyperref}}
\IfFileExists{xurl.sty}{\usepackage{xurl}}{} % add URL line breaks if available
\urlstyle{same}
\hypersetup{
  hidelinks,
  pdfcreator={LaTeX via pandoc}}

\author{}
\date{}

\begin{document}

\hypertarget{tidal-potential-measured-by-optical-clocks-on-a-coastal-inland-baseline}{%
\section{Tidal potential measured by optical clocks on a coastal-inland
baseline}\label{tidal-potential-measured-by-optical-clocks-on-a-coastal-inland-baseline}}

\begingroup\singlespacing

{[}Author list, order and affiliation assignments to be confirmed by the
authors{]}

\textsuperscript{1} Hefei National Laboratory, University of Science and
Technology of China, Hefei, China

\textsuperscript{2} Shanghai Research Center for Quantum Sciences and
CAS Center for Excellence in Quantum Information and Quantum Physics,
University of Science and Technology of China, Shanghai, China

\textsuperscript{3} Hefei National Research Center for Physical Sciences
at the Microscale and School of Physical Sciences, University of Science
and Technology of China, Hefei, China

\textsuperscript{4} Innovation Academy for Precision Measurement Science
and Technology, Chinese Academy of Sciences, Wuhan, China

\textsuperscript{5} Jinan Institute of Quantum Technology and CAS Center
for Excellence in Quantum Information and Quantum Physics, University of
Science and Technology of China, Jinan, China

\textsuperscript{6} Time and Frequency Geodesy Center (TFGC), School of
Geodesy and Geomatics, Wuhan University, Wuhan, China

\endgroup

\textbf{Tides alter the gravitational potential and thereby the relative
rates of distant clocks.\textsuperscript{1,2} Here we detect tidal
changes using a strontium clock in Shanghai and a ytterbium clock in
Wuhan, 692 km apart. The 545.9-hour record spans 353 days in 54
interrupted segments. Reducing data and prediction within each segment
retains the astronomical phase without joining their frequency levels.
The full tidal prediction has a fitted scale of \(0.96\pm0.11\) (one
statistical standard error). A joint fit with group-internal amplitude
ratios and phases fixed gives \(1.00\pm0.13\) for M₂+N₂ and
\(0.93\pm0.36\) for S₂+K₂. The lunar response is supported by the
observed amplitude, phase and spectral structure; the neighbouring group
is less precisely constrained. These measurements demonstrate tidal
potential sensing on a continental optical-clock baseline during
interrupted operation.}

\clearpage

Earth's gravitational potential changes as the Moon and Sun move and as
the oceans, atmosphere and solid Earth respond. Measuring these
variations provides information about mass redistribution and
deformation, and supports consistent height and time
references.\textsuperscript{3,4,5} Optical clocks make potential
differences observable through their relative rates: clocks at different
gravitational potentials accumulate time
differently.\textsuperscript{1,2} Fibre links connect clocks across
hundreds of kilometres and allow repeated measurements between the same
endpoints.\textsuperscript{6,7}

Laboratory redshift tests and remote clock comparisons have established
measurements of static potential differences.\textsuperscript{8,9,10,11}
Published chronometric-levelling experiments have also calculated the
tidal perturbation expected during a comparison.\textsuperscript{11}
Treating that perturbation as the measured signal tests a further
capability: can distant clocks detect the Earth's changing potential on
hourly timescales? The answer matters both for geophysical measurements
and for the time-dependent redshift corrections used in coordinated
clock comparisons.\textsuperscript{3,12,13}

Here we address this question with the strontium--ytterbium comparison
between coastal Shanghai and inland Wuhan. The observations are
interrupted, so the frequency level established during one run need not
be known in the next. The astronomical phase, however, remains defined
at every actual observation time. We combine the changes recorded within
separate segments to estimate the response to a complete tidal
prediction. We then ask how much of that response can be assigned to
neighbouring semidiurnal groups. The principal lunar tide M₂ repeats
every 12.42 h, whereas the principal solar tide S₂ repeats every 12 h.
Their changing relative phase makes observations on different dates
informative about their separate contributions.

\hypertarget{a-coastalinland-clock-baseline}{%
\subsubsection{A coastal--inland clock
baseline}\label{a-coastalinland-clock-baseline}}

The Shanghai and Wuhan clocks are 692 km apart and connected through
Hefei by a 1,350-km phase-stabilized fibre link (Fig. 1a). The apparatus
and frequency-comparison chain are described by Hou and
colleagues.\textsuperscript{14} Evaluations of the two clock systems and
the field-deployed transmission architecture provide the experimental
basis for the comparison.\textsuperscript{15,16,17} We take a positive
fractional-rate variation to mean Wuhan Yb running faster relative to
Shanghai Sr.~For displaying the measurements, the equivalent potential
is \(U=c^2y=-[\delta W_{\mathrm{Wuhan}}-\delta W_{\mathrm{Shanghai}}]\).

The coastal--inland geometry gives the two stations different responses
to ocean mass redistribution. The combined prediction contains the
direct astronomical forcing, the Earth's induced potential and clock
displacement, together with ocean-loading effects. We use the frozen
IERS-based solid prediction and FES2022 ocean elevations convolved with
a kernel that includes water attraction, induced potential and clock
displacement (Methods).\textsuperscript{18,19,20,21}

For M₂, the predicted solid and loading amplitudes are
\(3.43\times10^{-18}\) and \(1.42\times10^{-18}\) in fractional rate.
Their 19.8° phase difference gives a combined amplitude of
\(4.78\times10^{-18}\), equivalent to approximately
\(0.43\,\mathrm{m^2/s^2}\) (Fig. 1b). Loading is about 30\% of the
combined M₂ phasor amplitude. Its weighted root-mean-square variation is
33.1\% of the combined prediction on the sampled, segment-centred
record. These are model contributions; the experiment measures the
combined potential.

The record spans September 2025 to August 2026. Its 3,299 accepted
ten-minute means contain 545.9 nominal member-hours in 54 segments and
eight observing periods (Fig. 1c). The observations fall in eight
calendar months and on 43 UTC dates. Summer contributes 273.9 h, while
the earlier periods add different lunar--solar phases (Fig. 1d). The
analysis retains every accepted sample at its actual time.

\hypertarget{detecting-the-tidal-variation}{%
\subsubsection{Detecting the tidal
variation}\label{detecting-the-tidal-variation}}

Each segment is reduced about its own measured reference level. We
subtract the weighted mean separately from the observed rates and from
every prediction column. No fitted tidal amplitude enters this
operation. The remaining within-segment variations are then combined in
one estimate. The frequency levels are left independent across
interruptions, while the predicted phase fixes which part of the tidal
cycle each observation samples.

A single scale multiplying the complete tidal prediction gives
\(A_T=0.96\pm0.11\). The reported uncertainty follows the day-deletion
rule in Methods. Removing the tidal scale increases the weighted
residual sum of squares by 110.35. An eight-period restricted residual
sign-flip diagnostic gives \(P=0.0078\), at its finite-pattern
resolution limit (Methods).

The full record also reveals recurring modulation when ordered by M₂
phase (Fig. 2a). To display the lunar response, the observations and
unit lunar prediction undergo the same projection against segment
constants and the non-lunar groups. The displayed uncertainties
propagate that operation, and the support of each phase bin is retained.

Figure 2b,c returns to actual UTC observations. The displayed segment is
selected solely by its largest accepted member count: 42.5 h during
10--12 August 2026. Fixed UTC three-hour averages improve readability,
with the same averaging applied to the prediction. The complete
ten-minute measurements retain their full range below. The statistical
estimates use the original means from all 54 segments, rather than the
display averages or this selected interval.

Allowing the lunar response to deviate from the predicted phase gives
in-phase and quadrature coefficients of \(1.00\pm0.13\) and
\(-0.10\pm0.20\) (Fig. 3b). The astronomical phase is compatible with
the observed response. The phase fold, phase-relaxed fit and spectrum
describe the same record and contribute amplitude, phase and frequency
information.

\hypertarget{constraining-neighbouring-semidiurnal-responses}{%
\subsubsection{Constraining neighbouring semidiurnal
responses}\label{constraining-neighbouring-semidiurnal-responses}}

We jointly estimate the lunar group M₂+N₂, the neighbouring group S₂+K₂
and an auxiliary daily group. Relative amplitudes and phases are fixed
within each physical group. M₂ and N₂ are lunar constituents; S₂ is the
principal solar semidiurnal tide, while K₂ contains lunar and solar
forcing. Each semidiurnal group has its own coefficient in this common
model. The approximately daily constituents are retained to propagate
their covariance with the target responses.

The estimates are \(A_L=1.00\pm0.13\) and \(A_S=0.93\pm0.36\), relative
to the combined solid-Earth and loading prediction (Fig. 3a). Their
scaled least-squares correlation is −0.30. The weaker neighbouring
coefficient provides a less precise constraint, whose interpretation
depends more strongly on same-frequency disturbances. Conditional joint
regions and the targeted weak-group test are reported in Extended Data
Fig. 4 and Methods.

The observed semidiurnal spectrum has local maxima at 22.361 and 23.147
µHz (Fig. 3c). These lie near the known principal frequencies of M₂ and
S₂, whose astronomical spacing is approximately 0.78 µHz. The two
reported local maxima are 0.786 µHz apart. Spectra before and after
subtracting the fixed unit prediction retain the original gaps and use
the same segment projection. The potential-density axis is related to
the fractional-frequency density by \(S_W=c^4S_y\).

Single-line injections on the actual observation window produce about
8\% cross-power at the other principal frequency. Conditional scans
remove the competing principal sinusoid before examining each
neighbourhood (Extended Data Fig. 1). These calculations explain
sampling sidelobes and the remaining competition near S₂. The spectrum
displays frequency structure; the physical group amplitudes come from
the simultaneous template fit.

\hypertarget{measurements-across-observing-periods}{%
\subsubsection{Measurements across observing
periods}\label{measurements-across-observing-periods}}

Figure 4 shows the lunar and neighbouring responses estimated separately
in each of the eight observing periods. Each estimate uses only that
period's observations, with independent reference levels retained for
its constituent segments. The different durations and lunar--solar phase
coverage give unequal precision across periods. The full-record fit
combines these within-period variations at their known astronomical
phases.

A shared 24-h sine--cosine pair, motivated by daily operating schedules
and environmental variation, gives a lunar coefficient of
\(1.04\pm0.13\). Allowing a separate 12-h pair in every period gives
lunar \(0.75\pm0.24\) and neighbouring \(-5.10\pm15.01\). The broad
neighbouring interval reflects direct competition between the free 12-h
term and the astronomical template. These sensitivity scenarios,
including the shared 12-h fit, are reported in Extended Data Table 4.

The simultaneous round-trip transmission monitor covers 246.6
member-hours in summer. Adding it on the same observation times changes
the lunar coefficient by \(0.0061\pm0.0063\) and the neighbouring
coefficient by \(0.00594\pm0.01453\) (Extended Data Fig. 3). This
control constrains the monitored transmission route during the overlap.

The full-record response and lunar amplitude--phase agreement support
tidal-potential detection. The neighbouring coefficient has weaker
evidence under period-wise residual resampling, with a restricted wild
sign-flip diagnostic of \(P=0.0625\), and is reported as a component
constraint within the specified physical model.

The measurement extends a fixed-endpoint optical-clock comparison from a
static potential difference to its tidal variation. The segment-based
reduction uses the changes measured during each run and combines them at
their known astronomical phases. It provides a way to retain this
information through interruptions while keeping uncertain frequency
levels separate.

The substantial loading contribution makes this coastal--inland baseline
a useful test of the combined potential prediction. Separating loading
from the solid response will require observations that distinguish their
amplitude and phase more strongly. More continuous operation and
independent records of instrumental variation near 12 h would also
strengthen attribution of the neighbouring response. Related results
have been reported independently by a European optical-clock network in
a contemporaneous arXiv preprint.\textsuperscript{22}

\clearpage

\singlespacing\fontsize{11}{13}\selectfont

\hypertarget{references}{%
\subsection{References}\label{references}}

\begin{enumerate}
\def\labelenumi{\arabic{enumi}.}
\tightlist
\item
  Ludlow, A. D., Boyd, M. M., Ye, J., Peik, E. \& Schmidt, P. O. Optical
  atomic clocks. \emph{Rev.~Mod. Phys.} \textbf{87}, 637--701 (2015).
  \url{https://doi.org/10.1103/revmodphys.87.637}
\item
  Mehlstäubler, T. E., Grosche, G., Lisdat, C., Schmidt, P. O. \&
  Denker, H. Atomic clocks for geodesy. \emph{Rep.~Prog. Phys.}
  \textbf{81}, 064401 (2018).
  \url{https://doi.org/10.1088/1361-6633/aab409}
\item
  Voigt, C., Denker, H. \& Timmen, L. Time-variable gravity potential
  components for optical clock comparisons and the definition of
  international time scales. \emph{Metrologia} \textbf{53}, 1365--1383
  (2016). \url{https://doi.org/10.1088/0026-1394/53/6/1365}
\item
  Denker, H. et al.~Geodetic methods to determine the relativistic
  redshift at the level of \(10^{-18}\) in the context of international
  timescales: a review and practical results. \emph{J. Geod.}
  \textbf{92}, 487--516 (2018).
  \url{https://doi.org/10.1007/s00190-017-1075-1}
\item
  Schröder, S., Stellmer, S. \& Kusche, J. Potential and scientific
  requirements of optical clock networks for validating
  satellite-derived time-variable gravity data. \emph{Geophys. J. Int.}
  \textbf{226}, 764--779 (2021).
  \url{https://academic.oup.com/gji/article/226/2/764/6212226}
\item
  Lisdat, C. et al.~A clock network for geodesy and fundamental science.
  \emph{Nat. Commun.} \textbf{7}, 12443 (2016).
  \url{https://doi.org/10.1038/ncomms12443}
\item
  Riehle, F. Optical clock networks. \emph{Nat. Photon.} \textbf{11},
  25--31 (2017). \url{https://doi.org/10.1038/nphoton.2016.235}
\item
  Chou, C. W., Hume, D. B., Rosenband, T. \& Wineland, D. J. Optical
  Clocks and Relativity. \emph{Science} \textbf{329}, 1630--1633 (2010).
  \url{https://doi.org/10.1126/science.1192720}
\item
  McGrew, W. F. et al.~Atomic clock performance enabling geodesy below
  the centimetre level. \emph{Nature} \textbf{564}, 87--90 (2018).
  \url{https://doi.org/10.1038/s41586-018-0738-2}
\item
  Grotti, J. et al.~Geodesy and metrology with a transportable optical
  clock. \emph{Nat. Phys.} \textbf{14}, 437--441 (2018).
  \url{https://doi.org/10.1038/s41567-017-0042-3}
\item
  Takano, T. et al.~Geopotential measurements with synchronously linked
  optical lattice clocks. \emph{Nat. Photon.} \textbf{10}, 662--666
  (2016). \url{https://doi.org/10.1038/nphoton.2016.159}
\item
  Dimarcq, N. et al.~Roadmap towards the redefinition of the second.
  \emph{Metrologia} \textbf{61}, 012001 (2024).
  \url{https://doi.org/10.1088/1681-7575/ad17d2}
\item
  Lindvall, T. et al.~Coordinated international comparisons between
  optical clocks connected via fiber and satellite links. \emph{Optica}
  \textbf{12}, 843--852 (2025).
  \url{https://doi.org/10.1364/OPTICA.561754}
\item
  Hou, L. et al.~Yb/Sr optical clock comparison at \(2.2\times10^{-18}\)
  uncertainty over a 1350-km fiber link. \emph{arXiv}, 2610.34250v1
  (2026). \url{https://arxiv.org/abs/2610.34250v1}
\item
  Jia, Z.-P. et al.~Improved systematic evaluation of a strontium
  optical clock with uncertainty below \(1\times10^{-18}\).
  \emph{Metrologia} \textbf{63}, 025002 (2026).
  \url{https://doi.org/10.1088/1681-7575/ae449e}
\item
  Zhu, Q. et al.~Ytterbium lattice clock with systematic uncertainty of
  \(1.3\times 10^{-18}\) and instability at the \(10^{-19}\) level.
  \emph{arXiv}, 2606.10514v2 (2026).
  \url{https://arxiv.org/abs/2606.10514v2}
\item
  Chen, F.-X. et al.~10⁻²¹-level optical frequency dissemination over
  2067 km of noise-loaded field-deployed fiber network. \emph{Light Sci.
  Appl.} \textbf{15}, 276 (2026).
  \url{https://doi.org/10.1038/s41377-026-02299-1}
\item
  Petit, G. \& Luzum, B. (eds) IERS Conventions (2010). \emph{IERS
  Technical Note} \textbf{36} (2010).
  \url{https://iers-conventions.obspm.fr/content/tn36.pdf}
\item
  CNES. FES2022 (Finite Element Solution) Tidal model (Version 2024).
  CNES \url{https://doi.org/10.24400/527896/A01-2024.004} (2024).
\item
  Farrell, W. E. Deformation of the Earth by surface loads.
  \emph{Rev.~Geophys.} \textbf{10}, 761--797 (1972).
  \url{https://doi.org/10.1029/RG010i003p00761}
\item
  Han, D. \& Wahr, J. The viscoelastic relaxation of a realistically
  stratified earth, and a further analysis of postglacial rebound.
  \emph{Geophys. J. Int.} \textbf{120}, 287--311 (1995).
  \url{https://doi.org/10.1111/j.1365-246x.1995.tb01819.x}
\item
  Stahl, K. et al.~Optical clock network tracks tidal variations of the
  gravity potential. \emph{arXiv}, 2610.39635v1 (2026).
  \url{https://arxiv.org/abs/2610.39635v1}
\end{enumerate}

\clearpage

\singlespacing\fontsize{10}{12}\selectfont

\hypertarget{figure-legends}{%
\subsection{Figure legends}\label{figure-legends}}

\includegraphics{figures/figure1.pdf}

\textbf{Fig. 1 \textbar{} A coastal--inland baseline and its tidal
prediction.} \textbf{a,} Shanghai Sr and Wuhan Yb connected by the fibre
route through Hefei. The dashed line marks the 692-km baseline. Land
shading and contours illustrate the model M₂ clock-rate response in
\(10^{-18}\); sea contours show M₂ ocean-tide height in centimetres.
These are different physical quantities. \textbf{b,} Differential solid,
loading and combined M₂ phasors from the frozen numerical prediction,
with equal data-axis scales; the combined \(4.78\times10^{-18}\)
corresponds to about \(0.43\,\mathrm{m^2/s^2}\). \textbf{c,} Actual UTC
sampling and accepted member-hours in eight periods. \textbf{d,} S₂−M₂
phase at those same observations, including nodal corrections. P6--P8
are summer. The map fields and phasors are predictions; the lower panels
describe the measured observation window.

\clearpage

\includegraphics{figures/figure2.pdf}

\textbf{Fig. 2 \textbar{} Tidal variation in the measured clock rates.}
\textbf{a,} Twelve M₂-phase means of all 3,299 ten-minute observations.
The fixed minor remainder is removed, and data and unit lunar prediction
undergo identical weighted projection against segment levels and the
other physical groups. Error bars are one working standard error with
nuisance propagation; means are correlated. Row counts appear beneath
the phase labels. \textbf{b,} The segment with the largest accepted
member count, shown at real UTC. Points are fixed-UTC three-hour means
after independent weighted centring of data and prediction. Squares
apply the same averaging to the unit combined prediction; the line
evaluates it at the observed ten-minute times. Error bars propagate
centring. \textbf{c,} Every original ten-minute mean in that segment,
with its complete vertical range. No empirical daily sinusoid or
fit-dependent offset is removed. Left axes use \(U=c^2y=-\Delta W\);
right axes give fractional rate. Panels are views of the same
measurements, and display averaging does not change the primary sample
or inference.

\clearpage

\includegraphics{figures/figure3.pdf}

\textbf{Fig. 3 \textbar{} Amplitude, phase and frequency of the tidal
response.} \textbf{a,} The single complete-template response, followed
by the jointly fitted lunar and neighbouring semidiurnal coefficients.
Error bars are one statistical standard error, the larger of scaled LS
and UTC-day-deletion jackknife errors. The estimates share the same
observations and weights. \textbf{b,} Lunar in-phase and quadrature fit;
the cross marks \((1,0)\). Contours have covariance radii one and two
under scaled LS, not independently calibrated joint confidence levels.
\textbf{c,} The measured 21--24.5 µHz spectrum and the spectrum after
subtracting the fixed unit physical prediction. The scan retains all
3,299 original UTC rows and gaps, with equal weights and 54 segment
constants. Vertical markers are the prescribed M₂ and S₂ frequencies.
Left and right density axes obey \(S_W=c^4S_y\). Peaks, phase and
amplitudes use the same record and do not add independent detection
significance. Sampling cross-power and conditional scans are given in
Extended Data Fig. 1; joint group regions appear in Extended Data Fig.
4.

\clearpage

\includegraphics{figures/figure4.pdf}

\textbf{Fig. 4 \textbar{} Semidiurnal responses in eight individual
observing periods.} \textbf{a,} Lunar M₂+N₂ coefficients. \textbf{b,}
Neighbouring S₂+K₂ coefficients. Each point uses only the indicated
period, jointly estimates all three physical groups, and retains
independent reference levels for its original segments. All eight
periods are shown in calendar order, with 163, 203, 409, 447, 429, 601,
595 and 452 ten-minute means, respectively. Error bars show one
statistical standard error, the larger of scaled least squares and a
UTC-day-deletion jackknife computed within that period. The periods
contain 3, 3, 4, 5, 5, 9, 8 and 6 observed UTC dates. Both axes are
linear and show the full error ranges; their scales differ. Dashed lines
mark the unit prediction. The plotted periods use disjoint observations
and are fitted separately; the full-record estimate is obtained by the
common physical fit.

\clearpage

\fontsize{12}{14.4}\selectfont\doublespacing

\hypertarget{methods}{%
\subsection{Methods}\label{methods}}

\hypertarget{observations-frequency-sign-and-averaging}{%
\subsubsection{Observations, frequency sign and
averaging}\label{observations-frequency-sign-and-averaging}}

The Shanghai and Wuhan coordinates used by the frozen prediction are
31.130° N, 121.570° E and 30.529° N, 114.357° E. The measured ratio is
Wuhan Yb/Shanghai Sr.~With positive geodetic potential \(W\), its tidal
fractional variation is
\(-[\delta W_{\mathrm{Wuhan}}-\delta W_{\mathrm{Shanghai}}]/c^2\). The
period-relative beat conversion is
\(-d_p(f_{\mathrm{beat}}-f_{{\mathrm{ref}},p})/F_{1550}\), where
\(F_{1550}=1.933992\times10^{14}\) Hz and \(d_p=2\) in 2025 and 1 in
2026. Thus a 2026 beat deviation of 1 Hz has fractional magnitude
\(5.17\times10^{-15}\). The 2025 division factor is applied before
forming the small clock-rate variation.

The supplied valid-interval table defines 54 segments in 29 operating
groups and eight periods. Counter members are read across midnight
before binning. Beijing-time labels are converted to UTC. Summer group
corrections follow the decimal-arithmetic frequency-chain calculation;
early absolute correction entries remain incomplete. Their constant part
is represented by segment intercepts, while any time variation within a
segment remains a possible systematic contribution. Fitted intercepts
are never applied to join the runs for long-term stability estimation.

The membership audit contains 1,966,908 nominal one-second rows before
the minimum-bin-support cut. Three boundary rows reconcile its
difference from the printed interval total. Means are formed in nominal
600-s bins separately within each segment and require at least 180
retained members. The 3,299 accepted means contain 1,965,327 members, or
545.924 h. There are 3,291 unique bin centres: eight centres contain
separate partial means from different segments, with disjoint members.
The 353.39-day span contains observations on 43 UTC dates. The inherited
valid mask is retained; this work neither applies an additional
triangular filter nor selects points using tidal residuals. Exact
raw-file memberships and hashes are retained in the controlled archive.

\hypertarget{what-the-constituent-names-mean}{%
\subsubsection{What the constituent names
mean}\label{what-the-constituent-names-mean}}

Semidiurnal tides produce approximately two cycles per day. M₂ is the
principal lunar constituent; N₂ describes an additional lunar variation
associated with orbital ellipticity. S₂ is the principal solar
semidiurnal constituent, whereas K₂ contains lunar and solar forcing.
The names identify astronomical frequencies, not separate instruments or
independent experiments. The diurnal group contains K₁ (lunisolar), O₁
(principal lunar), P₁ (principal solar) and Q₁ (lunar elliptic). Their
periods are listed in Extended Data Table 1.\textsuperscript{23}
Internal relative amplitudes, phases and nodal modulation are fixed
within each physical group. A group response therefore tests that
complete specified pattern.

\hypertarget{combined-solid-and-loading-potential}{%
\subsubsection{Combined solid and loading
potential}\label{combined-solid-and-loading-potential}}

The physical target includes the direct astronomical potential, the
Earth's induced potential and the effect of clock displacement. The
selected frozen solid sequence implements these through the IERS radial
displacement \(u\) and a fixed effective-potential conversion,

\[y_{\mathrm{solid}}=-\frac{1+k_2-h_2}{h_2}\frac{g(u_{\mathrm{Wuhan}}-u_{\mathrm{Shanghai}})}{c^2},\]

with \(h_2=0.6078\), \(k_2=0.3010\) and \(g=9.794\,\mathrm{m/s^2}\).
This is a displacement-based approximation to the combined solid
potential. The same conversion also acts on degree-three and
frequency-dependent displacement contributions in the supplied sequence.
A previous implementation comparison found an approximately 18\% K₁
amplitude deficit relative to a potential-weighted reference. The daily
response is accordingly not used to infer a Love number. Uncertainty in
the solid model remains separate from the quoted statistical errors.

For loading, FES2022 sea-surface elevations \(\zeta\) are convolved with
the effective-potential kernel

\[\delta W_{\mathrm{load}}=\int \rho_w\zeta(\mathbf{r}',t)G_W(\psi)\,\mathrm{d}A,\qquad
G_W(\psi)=\frac{ga}{M_E}\sum_n(1+k'_n-h'_n)P_n(\cos\psi).\]

The three terms represent direct water attraction, induced potential and
clock displacement. PREM load Love numbers are retained through degree
696 with the archived high-degree continuation. Eighteen loading
constituents are synthesized with their astronomical nodal corrections.
The complex convention is \(f\,\mathrm{Re}[C\exp(-i\phi)]\); reversing
the exponential sign requires conjugating \(C\). The present analysis
reads the frozen constituent and time-series tables rather than
substituting an older loading sequence.

The eight named constituents form the three physical groups. Minor terms
remain once as a fixed remainder \(r_0\). On the actual weighted
segment-centred sampling, solid, loading and total root-mean-square
variations are respectively 2.928, 1.328 and \(4.011\times10^{-18}\).
Their ratio 0.331 is an amplitude-norm ratio, not an additive fraction
of total power. A separate loading coefficient is not estimated here;
the measured group responses test the combined prediction.

\hypertarget{segment-reduction-and-the-two-estimands}{%
\subsubsection{Segment reduction and the two
estimands}\label{segment-reduction-and-the-two-estimands}}

Let \(W_s\) be the original diagonal weight matrix in segment \(s\), and
define
\(\langle v\rangle_s=(\mathbf{1}^TW_sv)/(\mathbf{1}^TW_s\mathbf{1})\).
For each observation vector and each prediction column we form
\(\widetilde v_s=v_s-\mathbf{1}\langle v\rangle_s\). This operation uses
sample times and weights but no fitted tidal amplitude. For a physical
design matrix \(T_s\), the common coefficients solve

\[\widehat{\boldsymbol A}=\left(\sum_s\widetilde T_s^TW_s\widetilde T_s\right)^{-1}\sum_s\widetilde T_s^TW_s\widetilde y_s.\]

The stored segment matrices and vectors are sufficient statistics for
this estimate. All 54 segments enter the same fit; separately fitted
segment amplitudes are not averaged. An independent full-design QR fit
with 54 intercepts reproduces the three coefficients to within
\(3\times10^{-15}\). Recomputing the segment means after every deletion
also reproduces the original standard errors. No segment level is
transferred to create a continuous, calibrated long-term record.

For the full-response estimate, \(y\) is the observed fractional rate
and \(T=s_L+s_S+s_D+r_0\) is the complete frozen prediction. A single
coefficient \(A_T\) multiplies that entire column. The comparison is
between 54 segment constants alone and the same constants plus this one
template. For component estimation, we instead subtract the fixed minor
remainder once and fit

\[y_i-r_{0,i}=c_{s(i)}+A_Ls_{L,i}+A_Ss_{S,i}+A_Ds_{D,i}+\epsilon_i.\]

These estimands share the observations but answer different questions.
Both use \(W_{ii}=1/\sigma_i^2\) with the frozen Allan-based weights.
The period noise coefficients are estimated from within-segment
overlapping Allan deviations at 256--4,096 s after removal of the fixed
unit tide. No Allan window crosses a segment boundary. This does not use
stability calculated after joining fit-adjusted segment levels.

For Fig. 4, we repeat the same three-group estimation separately on each
of P1--P8. The data, fixed theoretical remainder and Allan-based weights
are inherited from the frozen primary design; each period retains its
original segment constants. The UTC-day jackknife is recomputed using
only the dates present in that period. All period fits and day-deletion
replicates retain full column rank and were independently checked
against an explicit segment-intercept QR calculation. The
individual-period errors use the same working standard-error rule as the
full record.

The daily group is an auxiliary physical component; its coefficient
remains \(0.58\pm0.54\) (Extended Data Table 2). The interval and
correction records do not establish calibrated zero transfer at the 46
boundaries where an eight-period intercept would be shared. They
document valid observing times but not signed boundary transfers with
uncertainties. The segment convention follows that experimental
information.

\hypertarget{statistical-errors-and-conditional-coverage}{%
\subsubsection{Statistical errors and conditional
coverage}\label{statistical-errors-and-conditional-coverage}}

For each group the reported standard error is the larger of
residual-scaled least squares and leave-one-UTC-day-out jackknife
errors. Deletions remove a calendar day, refit all remaining
coefficients and drop empty intercept columns. The jackknife variance is
\((G-1)G^{-1}\sum_g(\hat A_{(-g)}-\overline A)^2\). Period-deletion
errors are recorded separately. Extended Data Table 2 gives all
components and the same-data intercept comparisons. Observations are
repeated measurements of one clock pair, not thousands of independent
experiments.

Conditional simulations apply both designs to the same 3,299-row
sampling and heteroscedastic uncertainties. Each of four scenarios has
1,000 realizations: independent Gaussian noise; equal variance
contributions from Gaussian white noise and a 6-h Ornstein--Uhlenbeck
process within segments; the corresponding 24-h process within periods;
and standardized Student-\(t_5\) noise. The overall scale is the
observed segment-model residual scale; all true physical response
coefficients are one. Each realization recomputes both components of the
standard-error rule. Seed 2026110617 fixes the simulation. Coverage of
\(\hat A\pm1.96u\) is 95.1--97.4\% for the segment-model lunar
coefficient and 95.4--96.8\% for its neighbouring coefficient; Monte
Carlo standard errors are about 0.4--0.7 percentage points (Extended
Data Table 3). These are tests under declared noise laws, not a
calibration against every possible disturbance.

The new full-template coefficient uses the same standard-error rule. Its
scaled-LS, day-deletion, day-cluster CR1 and period-deletion standard
errors are respectively 0.0973, 0.1082, 0.1033 and 0.1216. The CR1
sandwich uses 43 UTC-day clusters and the finite-sample factor
\([G/(G-1)][(n-1)/(n-K)]\), with \(K=55\). The working interval
\(\widehat A_T\pm1.96u\) is 0.743--1.167. This is an explicitly labelled
working interval, rather than a confidence statement valid for every
possible noise process.

For targeted zero-response tests, we compare the weighted RSS of nested
models on the unchanged rows. The full-template test adds one entire
prediction to segment constants. The neighbouring-group test adds one
S₂+K₂ coefficient while retaining the lunar and daily groups and the
constants. Noise is generated on the true window with the original
heteroscedastic uncertainties and the observed full-model residual
scale. Each of the four declared noise scenarios has 20,000
realizations; seeds 2026110711 and 2026110709 fix the full-template and
neighbouring-group simulations. The exact residualized-template
projection computes the RSS improvement, and direct restricted/full
refits check that calculation. The full-template statistic is 110.353;
none of the 20,000 draws in any scenario reaches it. The plus-one tail
estimate is \(1/20,001\) in each scenario, not zero. For the
neighbouring statistic 10.943, the corresponding tail estimates are
0.00185, 0.12559, 0.01435 and 0.00205 in the listed noise-scenario
order.

Restricted period-wise residual sign flips enumerate all \(2^8=256\)
patterns, preserving each period's internal structure. The diagnostic
tail fractions are 2/256 for the full response and 16/256 for the added
neighbouring group. They rely on period independence and symmetry and
are not exact experimental randomization tests. The two-sided resolution
is 2/256. No Gaussian sigma value is inferred from either finite
enumeration.

The noise scenarios are conditional stress tests, not measured noise
laws. In separate 2,000-draw diagnostics per law, the observed
within-segment lag-one residual correlation is −0.148. The simulated
95\% ranges are −0.054--0.015 for Gaussian noise, 0.216--0.347 for the
6-h mixture, 0.072--0.202 for the 24-h mixture and −0.051--0.017 for
\(t_5\). Thus the positively correlated mixtures do not describe the
observed short-lag structure; neither do the other two laws fully
reproduce it. No scenario is selected because of its tidal tail
probability. No simulated record is added to the observations.

Extended Data Fig. 4 gives conditional two-group regions: the first
10,000 draws estimate error covariance and the empirical 95\%
Mahalanobis threshold; the other 10,000 assess coverage. Held-out joint
coverage is 94.3--94.8\%, with binomial Monte Carlo standard errors of
about 0.22 percentage points and additional calibration uncertainty.
These nominal regions describe the specified laws and do not replace the
reported marginal working errors with a universally calibrated joint
confidence region.

The phase-relaxed model adds the lunar quadrature while retaining the
prescribed M₂/N₂ relationship. Covariance contours use scaled least
squares. Subset and deletion fits refit the same model on their retained
rows. Shared-data estimates are not averaged. The simulation coverage
refers to two-sided intervals; the reported \(\pm\) values are standard
errors.

\hypertarget{phase-means-spectral-density-and-sampling-sidelobes}{%
\subsubsection{Phase means, spectral density and sampling
sidelobes}\label{phase-means-spectral-density-and-sampling-sidelobes}}

For the phase plot, \(N\) contains segment constants and the two
non-lunar groups. We project both \(y-r_0\) and the unit lunar
prediction with \(M_N=I-N(N^TWN)^{-1}N^TW\), and average in twelve fixed
M₂ phase bins. If \(B\) is weighted bin averaging, the mean covariance
is \(BM_NC_yM_N^TB^T\), where \(C_y\) is the scaled working observation
covariance. It retains correlations between means and propagates
nuisance estimation. The full unbinned response remains primary.

For the actual-time display, the representative segment maximizes
accepted nominal member count, with segment ID breaking ties. Segment 51
contains 255 ten-minute means and 153,000 members. Its 42.5 member-hours
extend from 10 August 2026 21:35 UTC to 12 August 15:55 UTC. Three-hour
display bins follow fixed UTC boundaries and remain inside the segment.
All 255 means are assigned; none is removed to improve the display.
Weighted averaging is identical for observed and unit-predicted centred
values. If a display bin has total weight \(w_b\) and its segment has
total weight \(w_s\), the centred-observation variance is
\(\widehat\sigma^2(1/w_b-1/w_s)\). This propagates the segment-mean
operation and creates correlations between display means. The curve at
ten-minute timestamps is a sampled physical prediction, not gap filling.
The potential display uses \(U=c^2y\) and the right axis retains the
fractional rate.

Figure 3c preserves the independently specified equal-weight,
54-intercept scan of the unchanged observed means, rather than replacing
it with the amplitude fit's weights. Every trial sine--cosine pair and
the data are projected against the same constants. For a reduction in
squared residuals \(\Delta\mathrm{RSS}(f)\) with dimensionless \(y\),
the nominal one-sided projection density is
\(S_y(f)=600\,\mathrm{s}\,\Delta\mathrm{RSS}(f)\) per Hz, or \(10^{-6}\)
times that value per µHz. Potential density is \(S_W=c^4S_y\); amplitude
density is \(c^2\sqrt{S_y}\). Uniform-grid checks reproduce the
conventional one-sided Fourier density. For interrupted data this
normalization is a stated projection density, not a Parseval-calibrated
variance over a filled year.\textsuperscript{24}

The frequency step is 0.002729 µHz; the inverse calendar span is
0.032752 µHz. Fine plotting steps do not improve physical resolution.
The listed peak positions are scan maxima; no constituent-frequency fit
or frequency-confidence interval is reported. M₂-only and S₂-only
astronomical injections use the same actual times, nodal factors and
equal-weight segment projection. Power at the opposite principal
frequency, divided by power at the injected frequency, gives 0.08027 and
0.08149. Conditional scans project out a freely phased sinusoid at the
other frequency before scanning the target neighbourhood. Amplitude
interference means that such corrections act on sampled waveforms or
complex coefficients, not by subtracting scalar power spectra. Gaps
alone do not predict the realization of measurement noise. Residual
density can describe remaining measurement fluctuations, but includes
possible model mismatch and cannot be called a pure instrumental noise
floor.

\hypertarget{instrument-and-daily-response-checks}{%
\subsubsection{Instrument and daily-response
checks}\label{instrument-and-daily-response-checks}}

The matched monitor calculation uses 1,482 rows with 887,882 nominal
members, or 246.634 h, in P6--P8. Both baseline and augmented fits have
identical members, segment labels and weights. Paired day-deletion
replicates propagate uncertainty in their difference. Extended Data Fig.
3 reports the response changes and dates. This comparison concerns the
monitored transmission route during the measured overlap.

On the selected data and model, a common 24-h sine--cosine pair gives
\(A_L=1.036\pm0.133\); separate 24-h pairs in each period give
\(1.003\pm0.136\). These daily sensitivity designs address a periodicity
motivated by temperature and operating schedules, without assigning an
independently measured thermal amplitude. A shared 12-h pair gives lunar
\(0.994\pm0.132\) and neighbouring \(2.931\pm1.011\). Separate pairs in
each period give \(0.751\pm0.242\) and \(-5.100\pm15.008\). The
normalized design condition number grows from 2.37 to 6.90 and 122.54,
respectively; all designs retain full numerical rank. These terms test
same-frequency identifiability and are not independently measured
disturbances. Extended Data Table 4 reports both groups and their
1.96-SE working intervals.

\hypertarget{ai-assisted-preparation}{%
\subsubsection{AI-assisted preparation}\label{ai-assisted-preparation}}

OpenAI Codex assisted English drafting, numerical replay and figure
preparation. {[}Authors to confirm their verification of the analyses,
citations and final text, and responsibility for the submitted work.{]}

\singlespacing\fontsize{11}{13}\selectfont

\hypertarget{data-availability}{%
\subsection{Data availability}\label{data-availability}}

The author review package contains the 3,299 derived means, fixed
templates, noise weights, figure source tables and immutable input
hashes. The original counters and membership ledger are held in the
controlled experimental archive. {[}Authors to supply an
editor-accessible location and confirm the permanent data-access
statement and repository identifier before submission.{]}

\singlespacing\fontsize{11}{13}\selectfont

\hypertarget{code-availability}{%
\subsection{Code availability}\label{code-availability}}

The package contains the independent QR replay, same-data design
comparison, conditional simulations, spectral calculations, covariance
propagation and figure scripts, with software versions and a
reproducible entry point. {[}Authors to confirm the code licence and
permanent repository identifier before submission.{]}

\singlespacing\fontsize{11}{13}\selectfont

\hypertarget{methods-references}{%
\subsection{Methods references}\label{methods-references}}

\begin{enumerate}
\def\labelenumi{\arabic{enumi}.}
\setcounter{enumi}{22}
\tightlist
\item
  NOAA. Harmonic constituents. \emph{NOAA Tides and Currents}.
  \url{https://tidesandcurrents.noaa.gov/harcon.html?id=8577330}
  (accessed 6 November 2026).
\item
  VanderPlas, J. T. Understanding the Lomb--Scargle periodogram.
  \emph{Astrophys. J. Suppl. Ser.} \textbf{236}, 16 (2018).
  \url{https://doi.org/10.3847/1538-4365/aab766}
\end{enumerate}

\clearpage

\singlespacing\fontsize{11}{13}\selectfont

\hypertarget{acknowledgements}{%
\subsection{Acknowledgements}\label{acknowledgements}}

{[}To be completed by the authors.{]}

\singlespacing\fontsize{11}{13}\selectfont

\hypertarget{funding-statement}{%
\subsection{Funding statement}\label{funding-statement}}

{[}Funding sources and grant identifiers to be completed by the
authors.{]}

\singlespacing\fontsize{11}{13}\selectfont

\hypertarget{author-contributions}{%
\subsection{Author contributions}\label{author-contributions}}

{[}To be completed and approved by all authors.{]}

\singlespacing\fontsize{11}{13}\selectfont

\hypertarget{competing-interests}{%
\subsection{Competing interests}\label{competing-interests}}

{[}Declaration to be supplied by the authors.{]}

\singlespacing\fontsize{11}{13}\selectfont

\hypertarget{additional-information}{%
\subsection{Additional information}\label{additional-information}}

\textbf{Correspondence:} {[}Corresponding author names and email
addresses to be supplied; mark them with an asterisk in the author
list.{]}

\clearpage

\singlespacing\sffamily\fontsize{7}{9}\selectfont

\hypertarget{extended-data}{%
\subsection{Extended Data}\label{extended-data}}

\includegraphics{figures/ed_figure2.pdf}

\textbf{Extended Data Fig. 1 \textbar{} Conditional semidiurnal scans.}
Equal-weight scans of the same 3,299 observations with 54 segment
intercepts. \textbf{a,} M₂ neighbourhood before and after projection of
a sine--cosine pair at S₂. \textbf{b,} S₂ neighbourhood before and after
projection at M₂. Dotted lines show the astronomical target frequencies.
These are descriptive projections on shared observations, not additional
detection probabilities.

\clearpage

\includegraphics{figures/ed_figure1.pdf}

\textbf{Extended Data Fig. 2 \textbar{} Same-data comparison of
reference-level designs.} \textbf{a,} Weighted segment-mean residuals
after the eight-period joint fit, ranging from \(-6.75\) to
\(12.00\times10^{-18}\). Segment constants make these means zero by
construction. \textbf{b,} Lunar point estimates after period deletion,
spanning 0.877--1.052 with eight period constants and 0.885--1.093 with
54 segment constants. The weighted residual sums of squares are 3,781.58
and 3,713.64; the latter has 46 additional intercept parameters.
Unweighted residual root-mean-square values are 24.114 and
\(23.934\times10^{-18}\), and reduced residual sums of squares are
1.1501 and 1.1455. Adjacent standardized residual correlations are
−0.127 and −0.148, using pairs within the same segment and separated by
at most 900 s. All comparisons retain identical observations, templates,
corrections and noise weights. These are residual diagnostics, not
measurements of physical jumps or zero-transfer calibration.

\clearpage

\includegraphics{figures/ed_figure3.pdf}

\textbf{Extended Data Fig. 3 \textbar{} Simultaneous summer transmission
control.} Changes caused by adding the fibre round-trip monitor to the
joint model, with paired UTC-day-jackknife standard errors. Baseline and
augmented fits use identical clock/monitor members, labels and weights:
1,482 means containing 887,882 nominal one-second members, or 246.634 h
in P6--P8, from 29 June to 25 August 2026. The lunar, neighbouring and
daily shifts are respectively \(0.00606\pm0.00627\),
\(0.00594\pm0.01453\) and \(0.00791\pm0.01388\). The paired monitor
constrains the measured transmission route on this summer overlap. The
main tidal estimate uses the entire 545.9-h clock record. The prior
transmission-architecture demonstration over 2,067 km is a separate
experiment.\textsuperscript{17}

\clearpage

\includegraphics{figures/extended4.pdf}

\textbf{Extended Data Fig. 4 \textbar{} Conditional uncertainty of the
neighbouring response.} \textbf{a,} Nominal 95\% joint regions for
\((A_L,A_S)\) under four declared noise laws on the original 3,299-row
design. The point is the observed estimate and the cross is \((1,1)\).
Each region uses 10,000 calibration and 10,000 validation simulations;
held-out coverage is 94.3--94.8\%. The regions are conditional on the
specified laws, none of which fully matches the measured short-lag
residual structure. \textbf{b,} Null tail curves for the added S₂+K₂
coefficient from 20,000 simulations per law. The vertical line marks the
observed weighted-RSS improvement; the dotted horizontal line is 0.05.
Simulations are separate from the measured data. The declared scenarios
provide sensitivity information, not four additional experiments or an
independently validated laboratory-noise model.

\clearpage

\textbf{Extended Data Table 1 \textbar{} Astronomical constituent
guide.} Names specify forcing and frequency; one response coefficient
multiplies each complete group with its internal ratios and phases
fixed.

\begin{longtable}[]{@{}lrll@{}}
\toprule\noalign{}
Symbol & Period\_h & Meaning & Group \\
\midrule\noalign{}
\endhead
\bottomrule\noalign{}
\endlastfoot
M2 & 12.421 & Principal lunar semidiurnal & Lunar \\
N2 & 12.658 & Lunar elliptic semidiurnal & Lunar \\
S2 & 12 & Principal solar semidiurnal & Near-12 h \\
K2 & 11.967 & Lunisolar semidiurnal & Near-12 h \\
K1 & 23.934 & Lunisolar diurnal & Diurnal \\
O1 & 25.819 & Principal lunar diurnal & Diurnal \\
P1 & 24.066 & Principal solar diurnal & Diurnal \\
Q1 & 26.868 & Lunar elliptic diurnal & Diurnal \\
\end{longtable}

\clearpage

\textbf{Extended Data Table 2 \textbar{} Estimator and uncertainty
comparison.} All directly replayed designs use the same 3,299 means,
physical templates and Allan-based weights. Errors are standard errors.

\begin{longtable}[]{@{}lllll@{}}
\toprule\noalign{}
Design & Lunar & Near\_12h & Diurnal & Status \\
\midrule\noalign{}
\endhead
\bottomrule\noalign{}
\endlastfoot
S54 & 1.000 ± 0.132 & 0.925 ± 0.365 & 0.579 ± 0.538 & Replayed \\
G29 & 0.984 ± 0.124 & 1.093 ± 0.298 & 0.759 ± 0.420 & Replayed \\
P8 & 0.987 ± 0.123 & 1.129 ± 0.288 & 0.678 ± 0.359 & Replayed \\
\end{longtable}

\begin{longtable}[]{@{}
  >{\raggedright\arraybackslash}p{(\columnwidth - 10\tabcolsep) * \real{0.1486}}
  >{\raggedleft\arraybackslash}p{(\columnwidth - 10\tabcolsep) * \real{0.1622}}
  >{\raggedleft\arraybackslash}p{(\columnwidth - 10\tabcolsep) * \real{0.1757}}
  >{\raggedleft\arraybackslash}p{(\columnwidth - 10\tabcolsep) * \real{0.1351}}
  >{\raggedleft\arraybackslash}p{(\columnwidth - 10\tabcolsep) * \real{0.1757}}
  >{\raggedleft\arraybackslash}p{(\columnwidth - 10\tabcolsep) * \real{0.2027}}@{}}
\toprule\noalign{}
\begin{minipage}[b]{\linewidth}\raggedright
Group
\end{minipage} & \begin{minipage}[b]{\linewidth}\raggedleft
Estimate
\end{minipage} & \begin{minipage}[b]{\linewidth}\raggedleft
Scaled\_LS
\end{minipage} & \begin{minipage}[b]{\linewidth}\raggedleft
Day\_JK
\end{minipage} & \begin{minipage}[b]{\linewidth}\raggedleft
Period\_JK
\end{minipage} & \begin{minipage}[b]{\linewidth}\raggedleft
Reported\_SE
\end{minipage} \\
\midrule\noalign{}
\endhead
\bottomrule\noalign{}
\endlastfoot
Lunar & 1.0004 & 0.132 & 0.1218 & 0.147 & 0.132 \\
Near-12 h & 0.9251 & 0.2993 & 0.3649 & 0.3531 & 0.3649 \\
Diurnal & 0.5792 & 0.3887 & 0.5381 & 0.7735 & 0.5381 \\
\end{longtable}

\clearpage

\textbf{Extended Data Table 3 \textbar{} Conditional 95\% interval
coverage.} Each scenario has 1,000 simulated realizations per design,
with fixed observed sampling and weights. Values are percentages for the
larger-of-LS-and-day-jackknife rule, using multiplier 1.96. They assess
the specified scenarios, not unknown systematic effects.

\begin{longtable}[]{@{}llrrr@{}}
\toprule\noalign{}
scenario & design & Diurnal & Lunar & Near12h \\
\midrule\noalign{}
\endhead
\bottomrule\noalign{}
\endlastfoot
50\% OU 24 h, period & P8 & 97.9 & 98 & 97.3 \\
50\% OU 24 h, period & S54 & 96.4 & 95.1 & 95.4 \\
50\% OU 6 h, segment & P8 & 94.7 & 95.7 & 96 \\
50\% OU 6 h, segment & S54 & 95.8 & 95.6 & 96.3 \\
Student t (df 5) & P8 & 96.8 & 96.3 & 96.2 \\
Student t (df 5) & S54 & 96.2 & 96.1 & 96.8 \\
Gaussian & P8 & 95.7 & 96.2 & 95.9 \\
Gaussian & S54 & 95.5 & 97.4 & 96.7 \\
\end{longtable}

\clearpage

\textbf{Extended Data Table 4 \textbar{} Same-data 12-h sensitivity.}
Each row retains the original 3,299 means, 54 segment levels,
astronomical group templates and Allan weights. Errors are the larger of
scaled LS and day-deletion jackknife standard errors. Working intervals
are estimate ±1.96 standard errors; they are not calibrated for all
correlated-noise or deterministic-disturbance models. A common 12-h
sine--cosine pair has one phase across the record; the period-wise
design gives each of eight periods its own pair. These fits assess
degeneracy, not the independently established magnitude of an
instrumental effect.

\begin{longtable}[]{@{}
  >{\raggedright\arraybackslash}p{(\columnwidth - 8\tabcolsep) * \real{0.1360}}
  >{\raggedright\arraybackslash}p{(\columnwidth - 8\tabcolsep) * \real{0.1680}}
  >{\raggedright\arraybackslash}p{(\columnwidth - 8\tabcolsep) * \real{0.2080}}
  >{\raggedright\arraybackslash}p{(\columnwidth - 8\tabcolsep) * \real{0.2240}}
  >{\raggedright\arraybackslash}p{(\columnwidth - 8\tabcolsep) * \real{0.2640}}@{}}
\toprule\noalign{}
\begin{minipage}[b]{\linewidth}\raggedright
Model
\end{minipage} & \begin{minipage}[b]{\linewidth}\raggedright
Lunar coefficient
\end{minipage} & \begin{minipage}[b]{\linewidth}\raggedright
Lunar working interval
\end{minipage} & \begin{minipage}[b]{\linewidth}\raggedright
Neighbouring coefficient
\end{minipage} & \begin{minipage}[b]{\linewidth}\raggedright
Neighbouring working interval
\end{minipage} \\
\midrule\noalign{}
\endhead
\bottomrule\noalign{}
\endlastfoot
Primary & 1.000 ± 0.132 & {[}0.742, 1.259{]} & 0.925 ± 0.365 & {[}0.210,
1.640{]} \\
Shared 12 h & 0.994 ± 0.132 & {[}0.735, 1.252{]} & 2.931 ± 1.011 &
{[}0.950, 4.913{]} \\
12 h per period & 0.751 ± 0.242 & {[}0.277, 1.224{]} & -5.100 ± 15.008 &
{[}-34.515, 24.316{]} \\
\end{longtable}

\end{document}
